\subsubsection{Différence de seuil et composition des couronnes}
\label{vi:cubo:mg07:umbral}

L'extension d'une cellule discrète conserve toutes les strates qui
apparaissent lorsqu'on augmente sa taille d'une unité. La différence de seuil
exprime cette opération sur le lecteur de cardinaux et permet de suivre
son comportement sous l'addition et la multiplication avant d'introduire une
réalisation métrique.

\begin{definition}[Différence de seuil]
Soit \(R\) un anneau commutatif unitaire et soit
\(f:\mathbb Z_{\geq0}\to R\). Nous définissons
\begin{equation}
 (D_\#f)(n)=f(n+1)-f(n).
 \label{vi:cubo:mg07:definicion}
\end{equation}
Pour un lecteur entier de cellules finies emboîtées, cette différence
enregistre le cardinal de la couronne incorporée à l'étape \(n\to n+1\).
\end{definition}

\begin{proposition}[Additivité et loi exacte du produit]
Pour \(f,g:\mathbb Z_{\geq0}\to R\), on a
\begin{align}
 D_\#(f+g)(n)&=D_\#f(n)+D_\#g(n),
 \label{vi:cubo:mg07:aditividad}\\
 D_\#(fg)(n)&=f(n+1)D_\#g(n)+g(n)D_\#f(n)
 \label{vi:cubo:mg07:producto-desplazado}\\
 &=f(n)D_\#g(n)+g(n)D_\#f(n)
   +D_\#f(n)D_\#g(n).
 \label{vi:cubo:mg07:producto}
\end{align}
\end{proposition}

\begin{proof}
La différence de la somme se sépare terme par terme. Pour le produit,
ajouter et soustraire \(f(n+1)g(n)\) donne
\[
 \begin{split}
 f(n+1)g(n+1)-f(n)g(n)
 &=f(n+1)\bigl(g(n+1)-g(n)\bigr)\\
 &\quad+g(n)\bigl(f(n+1)-f(n)\bigr).
 \end{split}
\]
Cela prouve \eqref{vi:cubo:mg07:producto-desplazado}. La substitution
\(f(n+1)=f(n)+D_\#f(n)\) produit
\eqref{vi:cubo:mg07:producto} avec son terme bilinéaire complet.
\end{proof}

Si \(A_n\subseteq A_{n+1}\) et \(B_n\subseteq B_{n+1}\) sont des
ensembles finis, la couronne du produit se décompose en trois parties
disjointes :
\[
 \begin{split}
 (A_{n+1}\times B_{n+1})\setminus(A_n\times B_n)
 &=\bigl(A_n\times(B_{n+1}\setminus B_n)\bigr)\\
 &\quad\sqcup\bigl((A_{n+1}\setminus A_n)\times B_n\bigr)\\
 &\quad\sqcup\bigl((A_{n+1}\setminus A_n)
                 \times(B_{n+1}\setminus B_n)\bigr).
 \end{split}
\]
Prendre les cardinaux donne exactement \eqref{vi:cubo:mg07:producto}. Le terme
bilinéaire compte les positions nouvelles dans les deux facteurs ; son lieu de résidence
est l'intersection des deux extensions, avec l'incidence conservée.

\begin{proposition}[Loi binomiale de la couronne dimensionnelle]
Pour \(d\geq1\), soit \(f_d(n)=n^d\). Alors
\begin{equation}
 D_\#f_d(n)=(n+1)^d-n^d
 =\sum_{k=0}^{d-1}\binom dk n^k
 =\sum_{r=1}^{d}\binom dr n^{d-r}.
 \label{vi:cubo:mg07:binomio}
\end{equation}
\end{proposition}

\begin{proof}
Le développement binomial de \((n+1)^d\) contient \(n^d\) et les \(d\) termes restants de la première somme. La substitution \(r=d-k\) donne la seconde. Sur le support \(\{1,\ldots,n+1\}^d\), chaque position ajoutée possède un ensemble non vide de coordonnées égales à \(n+1\). Si cet ensemble contient \(r\) éléments, il y a \(\binom dr\) choix de positions et \(n^{d-r}\) choix des coordonnées précédentes. Ces strates sont disjointes et couvrent la couronne ; elles prouvent donc également la formule par incidence finie.
\end{proof}

Pour \(d=3\), les trois strates ont les cardinaux \(3n^2\),
\(3n\) et \(1\). Au passage nonadique--décimal, \(n=9\), l'
opération donne \(243+27+1=271\) ; l'exclusion du sommet final laisse
\(270\) positions. La complétion cubique qui suit explicite l'
union de cette couronne avec les \(729\) positions antérieures. La loi
de seuil conserve ainsi tant le total que la répartition selon le nombre
de nouvelles coordonnées. Les réalisations réticulaires ultérieures
transportent leurs incidences au moyen des applications spécifiques de chaque
construction.
